\documentclass{article}
\usepackage[utf8]{inputenc}
\usepackage[T1]{fontenc}
\usepackage[a4paper,total={6in,8in}]{geometry}
\usepackage{amsmath,amssymb,amsfonts}
\usepackage{graphicx}
\usepackage{hyperref}
\usepackage{enumitem}
\usepackage{listings}
\usepackage{xcolor}
\DeclareMathOperator{\Var}{Var}
\DeclareMathOperator{\Cov}{Cov}
\providecommand{\tightlist}{\setlength{\itemsep}{0pt}\setlength{\parskip}{0pt}}
\lstset{basicstyle=\ttfamily\scriptsize,breaklines=true,columns=fullflexible,
  keepspaces=true,showstringspaces=false,frame=single}
\hypersetup{colorlinks=true,urlcolor=blue,linkcolor=blue}
\setlength{\emergencystretch}{2em}

\title{Math 126 Problem 3: Count-min Works Super Well}
\author{} % Put your name here.
\date{Spring 2027}
\begin{document}
\maketitle

\hypertarget{submission-instructions}{%
\subsection*{Submission Instructions}\label{submission-instructions}}

Please upload your work to Gradescope by \textbf{11:59pm on Monday, February 1, 2027 (Pacific time).}

\begin{itemize}
\item
  You are encouraged to discuss ideas and work with your classmates. However, you \emph{must acknowledge} your collaborators at the top of each solution on which you collaborated with others and you \emph{must write} your solutions independently.
\item
  Your solutions to theory questions must be written legibly, or typeset in LaTeX or markdown.
\item
  I recommend that you write your solutions in a Python notebook using Google Colab.
\item
  You should submit your solutions as a \emph{single PDF} via the assignment on Gradescope.
\end{itemize}

\textbf{Grading:} The point of the problem set is for \emph{you} to learn. To this end, I hope to disincentivize the use of LLMs by \emph{not} grading your work for correctness. Instead, you will grade your own work by comparing it to my solutions. This self-grade is due the Friday \emph{after} the problem set is due, also on Gradescope.

\noindent Original notebook: \href{https://colab.research.google.com/drive/1qu-ejYqkDLDEhy8UQOnguYA1U2qA9G_R?usp=sharing}{Google Colab}.

% Source notebook cell 2
\hypertarget{count-min-works-super-well}{%
\subsection*{Count-min Works Super Well}\label{count-min-works-super-well}}

We showed that Count-Min can estimate the frequency of any item in a stream of \(n\) items up to additive error \(\frac{1}{m} n\) using \(O(m)\) space. In practice it is often observed that this bound is pessimistic: the algorithm performs better than expected. In this problem, you will establish one reason why.

For any positive integer \(m\), let \(f_1, \ldots, f_{m}\) be the frequencies of the \(m\) most frequent items in our stream. Let \(C = n - \sum_{i=1}^{m} f_i\). In general, we can have that \(C \ll n\). For example, it has been observed that up to \(95\%\) of YouTube video views come from just \(1\%\) of videos. Prove that using \(O(m)\) space, Count-Min actually returns an estimate \(\tilde{f}(v)\) to \(f(v)\) for any item \(v\) satisfying: \begin{align*}
    f(v) \leq \tilde{f}(v) \leq f(v) + \frac{2}{m} C
\end{align*} with \(9/10\) probability. This is strictly better than the \(\frac{1}{m} n\) error bound shown in class.

\hypertarget{part-a}{%
\subsection*{Part A}\label{part-a}}

Explain why it suffices to show the equation holds with any constant probability \(c > 0\).

Then write the estimate \(\tilde{f}(v) = A[h(v)]\) as the true frequency \(f(v)\) and two error terms; the first error term is for the \(m\) most frequent items that could collide with \(v\) and the second error term is for the remaining items that could collide with \(v\)..

\hypertarget{part-b}{%
\subsection*{Part B}\label{part-b}}

Prove that the first error term is 0 with constant probability; that is, there is a constant probability that none of the \(m\) most frequent items collide with \(v\).

Hint: Use one of these inequalities: \begin{align*}\frac1{e} \geq \left(1-\frac1{m}\right)^m \geq \frac{1}{2e}.\end{align*} You can check these inequalities for yourself on \href{https://www.desmos.com/calculator/bok5upfowo}{Desmos}.

Following the analysis in class, prove that the second error term is at most \(\frac{2}{m}C\) with constant probability.

Show how both results together imply that the equation holds with a constant probability.

\end{document}
